\begin{abstract}
Remote-sensing change detection requires spatially coherent change masks and accurate boundary delineation. Existing CNN and transformer-based Siamese architectures model temporal differences effectively, but region-level change extent and boundary-level contour evidence are typically fused into a single representation, which can limit recovery of fine object boundaries. This paper presents \textbf{\ours}, a MambaVision-based binary change detection framework that explicitly decouples temporal evidence into complementary region-aware and boundary-aware responses. A shared MambaVision encoder extracts aligned multi-scale features from bi-temporal image pairs. The proposed \textbf{Differential Region--Boundary Interaction} (\drbi) module applies separate bounded gates to a temporal feature stream constructed from both absolute and signed bi-temporal differences, producing distinct region and boundary pathways before multi-scale aggregation. An adaptive receptive-field decoder aggregates multi-scale context through parallel dilated convolution branches with per-image learned attention, and a bounded residual refinement head adjusts coarse predictions near boundaries with a controlled correction magnitude. On DSIFN-CD, averaged over two canonical runs, the method achieves F1\,=\,95.67 and IoU\,=\,91.71. On WHU-CD, it achieves F1\,=\,95.15 and IoU\,=\,90.76, using 65.40M parameters. A controlled ablation over seven variants (A0--A6) isolates the contribution of each component.
\end{abstract}
